\documentclass[10pt,a4paper]{amsart}
\usepackage[margin=19mm]{geometry}
\usepackage{amsmath,amssymb,mathtools}
\usepackage[scr=rsfs,cal=euler]{mathalfa}
\usepackage{xfrac}
\usepackage[unicode,hidelinks,bookmarks=false]{hyperref}
\newcommand{\ie}{i.e.\ }
\newcommand{\cf}{cf.\ }
\newcommand{\resp}{respectively\ }
\newcommand{\Subsection}{\subsection}
\providecommand{\nobreakdash}{\nobreak\hbox{-}}
\setlength{\emergencystretch}{2em}
\multlinegap0pt
\usepackage{bm}
\usepackage{bbm}
\usepackage{dsfont}
\usepackage{xspace}
\usepackage{cancel}
\usepackage{mathtools}
\usepackage[shortlabels]{enumitem}
\usepackage{amsthm}
\makeatletter
\@ifundefined{theorem}{\newtheorem{theorem}{Theorem}}{}
\@ifundefined{lemma}{\newtheorem{lemma}[theorem]{Lemma}}{}
\@ifundefined{proposition}{\newtheorem{proposition}[theorem]{Proposition}}{}
\makeatother
\makeatletter
\newcommand{\lla}{\langle\!\langle}
\expandafter\ifx\csname maincorollary\endcsname\relax
\newenvironment{maincorollary}[1]
  {\renewcommand\theinnermaincorollary{\ref{#1}}\innermaincorollary}
\fi
\expandafter\ifx\csname maintheorem\endcsname\relax
\newenvironment{maintheorem}[1]
  {\renewcommand\theinnermaintheorem{\ref{#1}}\innermaintheorem}
\fi
\newcommand{\rra}{\rangle\!\rangle}
\makeatother
\title[On canonical roots of fractional ideals]{On canonical roots of fractional ideals}
\author{Daniel M.H. van Gent}
\date{Source: arXiv:2607.19135v1 (CC BY 4.0)}
\begin{document}
\maketitle

\noindent\emph{Source and licence.} On canonical roots of fractional ideals. Version arXiv:2607.19135v1 (CC BY 4.0), distributed under the Creative Commons Attribution 4.0 International licence, as recorded in the arXiv licence metadata for this exact version and shown on the abstract page https://arxiv.org/abs/2607.19135v1. Full rights notice: licenses/arxiv-2607.19135-CC-BY-4.0.txt.\par\smallskip

\noindent\textit{Editorial scope.} Counted unit: the Proposition whose source label is \texttt{prop:closure} in the article (source0.tex lines 747--754, source label \texttt{prop:closure}), delivered together with the article's own complete original proof of it (source0.tex lines 755--782, one \texttt{proof} environment of 28 lines). No mathematics is written, repaired or re-derived: statement and proof are the article's own text line for line. Required context retained uncounted: lem:closure\_is\_monoid (lines 713--715); lem:coprime\_basis\_partial\_order (lines 735--737). Every definition, display and label of those lines is retained unchanged. Omitted from this package and not redistributed: 1--712, 716--734, 738--746, 783--1013. Nothing inside a retained excerpt is omitted or altered: every retained body is delivered exactly as the article's lines are written, so no omission receipt of any kind accompanies this package. Citations: the retained excerpts carry no citation command at all, so no bibliography entry of the article is transcribed and no reference is redistributed. Reference adaptation: none was needed and none was made; every reference inside the delivered unit resolves inside it, so no unresolved reference or citation is produced. Printed slips kept literal: no printed word, symbol, hypothesis or number of the counted statement or of its proof was altered. Layout and class: the article's own class is \texttt{amsart} with packages including inputenc, babel, geometry, amsmath, amsthm, amscd, amssymb, mathtools, verbatim, wrapfig, bbm, cite, hyperref, tikz-cd; the wrapper is the lane's pinned \texttt{amsart} editorial wrapper at 10pt on A4, which restates the theorem-like environments the retained text needs and adds the article's own macro definitions that the retained text needs (\texttt{\textbackslash lla}, \texttt{\textbackslash rra}), with extra wrapper packages (bm, bbm, dsfont, xspace, cancel, mathtools, enumitem). The delivered numbering is the unit's own. Assets: no retained span contains a figure environment, a graphics-inclusion command, a PDF-page inclusion, a table or a TikZ picture; no image file is copied into this package, the package declares no asset dependency and no image file is redistributed with it. Licence: version arXiv:2607.19135v1 of this article is distributed under the Creative Commons Attribution 4.0 International licence, as recorded in the arXiv licence metadata for this exact version and shown on the abstract page https://arxiv.org/abs/2607.19135v1. A full rights notice is delivered at licenses/arxiv-2607.19135-CC-BY-4.0.txt. Characters: every retained excerpt is pure ASCII, so the delivered engine, XeLaTeX with its default Latin Modern text font, reports no missing character.
\begin{lemma}\label{lem:closure_is_monoid}
Let \(R\) be a commutative ring and \(C\) a set of invertible ideals contained in \(R\) which are pairwise coprime. Then \(\textup{cl}(C)=\lla C\rra\).
\end{lemma}

\begin{lemma}\label{lem:coprime_basis_partial_order}
For a commutative ring \(R\) and a set \(X\) of fractional ideals contained in \(R\) we may equip the set of coprime bases of \(X\) with a partial order where \(C\leq D\) if and only if \(\lla C\rra\subseteq\lla D\rra\).
\end{lemma}

\begin{proposition}\label{prop:closure}
Let \(R\) be a Noetherian commutative ring and \(X\) a set of fractional ideals contained in \(R\). Then:
\begin{enumerate}[nosep,label=\textup{(\roman*)}]
\item \(X\) has a coprime basis if and only if \(\textup{cl}(X)\subseteq\mathcal{I}(R)\);
\item if \(X\) has a coprime basis, then it has a unique minimal one;
\item if \(C\) is a coprime basis of \(X\), then \(C\) is minimal if and only if \(C\subseteq\textup{cl}(X)\).
\end{enumerate}
\end{proposition}

\begin{proof}
(i) Suppose \(X\) has a coprime basis \(D\).
Then \(X\subseteq \lla D\rra = \textup{cl}(D)\) by Lemma~\ref{lem:closure_is_monoid}, so \(\textup{cl}(X)\subseteq\textup{cl}(D)=\lla D\rra \subseteq\mathcal{I}(R)\). 
Suppose instead that \(\textup{cl}(X)\subseteq \mathcal{I}(R)\).
We will show that
\[C=\{ \mathfrak{a}\in \textup{cl}(X) \,|\, \forall\, \mathfrak{b}\in \textup{cl}(X),\, \mathfrak{a}\subsetneq \mathfrak{b} \Leftrightarrow \mathfrak{b}=R\}\]
is a coprime basis of \(X\).

First, note that the elements of \(C\) are pairwise coprime: For \(\mathfrak{a},\mathfrak{b}\in C\) we have \(\mathfrak{a}+\mathfrak{b}\in \textup{cl}(X)\). If \(\mathfrak{a}\subsetneq \mathfrak{a}+\mathfrak{b}\), then \(\mathfrak{a}+\mathfrak{b}=R\) by definition of \(C\), and similarly when \(\mathfrak{b}\subsetneq\mathfrak{a}+\mathfrak{b}\). Otherwise \(\mathfrak{a}=\mathfrak{a}+\mathfrak{b}=\mathfrak{b}\).
Second, we show \(\textup{cl}(X)\subseteq\lla C\rra\) using Noetherian induction:
Certainly \(R\in\lla C\rra\). 
Now let \(\mathfrak{a}\in \textup{cl}(X)\setminus\{R\}\) and suppose \(\mathfrak{c}\in\lla C\rra\) for all \(\mathfrak{c}\in \textup{cl}(X)\) with \(\mathfrak{a}\subsetneq \mathfrak{c}\).
Either \(\mathfrak{a}\in C\), or there is some \(\mathfrak{b}\in \textup{cl}(X)\) such that \(\mathfrak{a}\subsetneq \mathfrak{b} \subsetneq R\), in which case \(\mathfrak{b}, (\mathfrak{a}:\mathfrak{b})\in\lla C\rra\) by the induction hypothesis and hence \(\mathfrak{a}\in\lla C\rra\). Thus \(C\) is a coprime basis for \(X\), as was to be shown.

(ii) Suppose now that \(X\) has a coprime basis. 
We will show that \(C\) as in (i) is the unique minimal coprime basis.
Let \(D\) be any coprime basis of \(X\).
We have \(C\subseteq\textup{cl}(X)\), so \(\textup{cl}(C) \subseteq \textup{cl}(X)\).
On the other hand, \(X\subseteq \textup{cl}(C)\), so \(\textup{cl}(C)=\textup{cl}(X)\).
Similarly for \(D\) we have \(\textup{cl}(X)\subseteq\textup{cl}(D)\).
Hence \(\lla C \rra =\textup{cl}(C)=\textup{cl}(X)\subseteq\textup{cl}(D)=\lla D\rra\) by Lemma~\ref{lem:closure_is_monoid}.
Thus \(C\leq D\), as was to be shown.

(iii) It is clear that the minimal coprime basis from (ii) satisfies \(C\subseteq\textup{cl}(X)\).
Let \(D\) be any coprime basis of \(X\) such that \(D\subseteq\textup{cl}(X)\).
Then as before we obtain \(\lla D\rra = \textup{cl}(D)=\textup{cl}(X)\).
Hence \(\lla C \rra = \textup{cl}(X)=\lla D\rra\) and \(C=D\) by Lemma~\ref{lem:coprime_basis_partial_order}. Hence \(D\) is minimal.
\end{proof}

\end{document}
